\subsection*{The actual quintic resolvent homomorphisms}

\statement{B.53}{Proved}{The two resolvents and the corrected third target}
Let $q_5:\widehat B\twoheadrightarrow A_5$, let
$T=A_4\leq A_5$ be a point stabilizer in its five-letter action,
and put $J=q_5^{-1}(T)=\widehat{\Gamma_T}$. Geometrically mark
$\Gamma_T=C_2*C_3*C_3$ by the quintic of B.11. Its orbifold is
\[
 \mathscr X=\bigl(\mathbb P^1_t;\ 3\text{ of order }2,
       \ r_+,r_-\text{ of order }3;\ \infty\text{ a cusp}\bigr),
 \qquad r_\pm=\frac{-5\pm3\sqrt{-15}}2.
\]
Here a marked orbifold means the root stack at the finite cone
points, with the cusp removed; its geometric fundamental group
is the full profinite completion of the displayed discrete group.
Set
\[
 k=\mathbb Q(\sqrt{-15}),\qquad
 D=\{\Phi\in\widehat{GT}_0:
                  \chi_{-15}(\lambda(\Phi))=1\}.
\]
The permutation of the two order-three inertia classes under
restriction of $\Phi$ is exactly $\chi_{-15}(\lambda(\Phi))$.
Filling the cone point $r_-$ or $r_+$ therefore gives continuous
homomorphisms of profinite groups
\[
             \mathfrak T_+,\mathfrak T_-:
                         D\longrightarrow\widehat{GT}_0.
                                                               \tag{B.53.1}
\]
Both preserve $\lambda$, both are the identity on the Galois
image of $G_k$, and
$\mathcal K\subseteq\ker\mathfrak T_+\cap\ker\mathfrak T_-$.
The target identifications in (B.53.1) are the geometric ones
specified below; they are part of the assertion.

Filling the order-two point instead gives a continuous homomorphism
\[
 \mathfrak T_2:D\longrightarrow
  \operatorname{Out}_{\partial}(C_3\amalg C_3),
                    \qquad \mathcal K\subseteq\ker\mathfrak T_2.
                                                               \tag{B.53.2}
\]
The target preserves each named factor class up to the common
unit power and preserves the cusp class up to that power.
It is the outer group of the weighted three-point orbifold
$(\mathbb P^1;0,1\text{ of order }3;\infty\text{ a cusp})$.
It is not identified with $\widehat{GT}_0$ in this statement.
Thus the first two maps of P4 are established, while its literal
third-target identification requires the correction in (B.53.2).

\begin{proof}
We give the finite permutation computation first, since it fixes
the domain on the whole gentle group, independently of arithmetic
examples. Every automorphism of $\widehat B$ preserves
$\widehat K$ by B.4. Its induced automorphism of $A_5$ carries
$T$ to a conjugate, since the five subgroups $A_4$ form one
conjugacy class. Moreover
\[
                    N_{\widehat B}(J)=J:
\]
apply $q_5$ and use $N_{A_5}(T)=T$. Hence outer restriction to
$J$ is canonical. Explicitly, choose an inner adjustment carrying
the image of $J$ back to $J$ and restrict; two choices differ by
an element of its normalizer $J$, and thus give the same outer
class. This also proves that restriction is a homomorphism.
It is continuous: the required conjugating coset is chosen from
the finite five-coset set, and restriction is continuous on each
of these finitely many open pieces.

Choose such an adjusted representative $\Phi$. Its quotient
on $A_5$ is conjugation by a unique $\nu\in S_5$, with $\nu$
fixing the letter stabilized by $T$. Write
\[
          \Phi(r)=A r^\lambda A^{-1},\qquad
          w=q_5(A)^{-1}\nu\in S_5.
\]
The cusp $q_5(\tau)$ is a five-cycle. Its two conjugacy classes
in $A_5$ correspond to square and nonsquare unit exponents
modulo five. Thus
\[
                \operatorname{sgn}(\nu)=\left(\frac\lambda5\right).
                                                               \tag{B.53.3}
\]
The displayed formula for $\Phi(r)$ gives
$wq_5(r)w^{-1}=q_5(r)^\lambda$. The normalizer in $S_5$ of
this three-cycle preserves its three-letter support and its
two fixed letters. On the support the sign is
$(\lambda/3)$, since inversion of a three-cycle requires an
odd support permutation. On all five letters the sign is
$(\lambda/5)$, because $q_5(A)$ is even. The sign of the
permutation of its two fixed letters is consequently
\[
                 \left(\frac\lambda5\right)
                 \left(\frac\lambda3\right)
                           =\chi_{-15}(\lambda).           \tag{B.53.4}
\]

These two fixed letters are exactly the two conjugacy classes
of order-three factors of $\Gamma_T$. To check the transport,
a class represented by $g^{-1}rg\in\Gamma_T$ corresponds to
the fixed letter $q_5(g)o$, where $o$ is the $T$-fixed letter.
Its image has representative
$(A^{-1}\Phi(g))^{-1}r^\lambda(A^{-1}\Phi(g))$, whose fixed
letter is $wq_5(g)o$. This proves (B.53.4) for the actual
restriction, rather than for an unrelated finite permutation.
The unique order-two class is preserved. The unique cusp is
$\langle\tau^5\rangle$. A conjugator in $\widehat B$ taking
its image to $\tau^{5\lambda}$ can be multiplied by a power
of $\tau$ to lie in $J$, because $A_5=T\langle q_5(\tau)\rangle$.
This adjustment does not change the conjugated cusp element.
Thus its exponent after restriction is still $\lambda$.

For $\Phi\in D$, each of the two closed normal closures of
the order-three factors is preserved. Killing either factor
therefore induces an automorphism of $C_2\amalg C_3$,
with the two torsion classes and the cusp class carried to
their $\lambda$-powers. This construction respects composition
and inner classes. The Harbater--Schneps comparison in A.1,
or its centerless-extension form in B.5 applied to the
characteristic level-two kernel, identifies this peripheral
outer group with $\widehat{GT}_0$. In that comparison one
adjusts the action on $S_3$ to the identity; the adjustment
is unique up to the level-two kernel since $Z(S_3)=1$.
It therefore does not introduce an additional quotient by
geometric $S_3$. This proves the types and multiplier assertion
of (B.53.1). Killing the order-two factor instead leaves
$C_3\amalg C_3$ and proves (B.53.2) by the same quotient argument.

We now specify the geometric markings that make the arithmetic
statement exact. After filling $r_-$, the affine function
\[
                j'_+(t)=1728\frac{t-r_+}{3-r_+}
\]
sends the remaining order-three point to zero, the order-two
point to $1728$, and the cusp to infinity. The corresponding
formula with $r_-$ defines $j'_-$. Both are defined over $k$.
They identify the filled root stacks with the standard
$(2,3,\infty)$ modular orbifold over $k$. With geometric
fundamental groups identified by these maps, functoriality
of the arithmetic exact sequences gives
\[
           \mathfrak T_+(\sigma)=\sigma
                =\mathfrak T_-(\sigma)\quad(\sigma\in G_k).
                                                               \tag{B.53.5}
\]
Changing a base-point path introduces an inner automorphism,
already removed in these targets. An arbitrary abstract
peripheral identification could instead conjugate the
arithmetic embedding; it is not the identification chosen
here. For (B.53.2) use
$t'=(t-r_+)/(r_--r_+)$ to put its two surviving cone points
at zero and one over the same field. Its restriction on
$G_k$ is the natural arithmetic outer action on that
weighted three-point orbifold.

Finally let $h\in\mathcal K$. By B.30 it is an actual
limit of mapping classes centralizing $T$. The completed
covering-normalizer sequence descends these classes to the
four-special-point quotient orbifold $\mathscr X$.
Their outer action on $J$ agrees with the restriction just
constructed: both extend the same outer action on
$\widehat K$, induce the identity on $T$, and the
centerless extension comparison makes them equal.
The multiplier-one condition gives (B.53.4) equal to one,
so both order-three points are individually fixed; the
order-two point and cusp are unique and fixed as well.
After any of the three fillings the descended mapping
class is a pure mapping class of a three-marked sphere,
and hence is the identity. Taking completions preserves
this zero homomorphism. Thus all three specified images
of $h$ are trivial, as asserted.
\end{proof}

B.53 settles the construction and the arithmetic comparison of
these resolvents. B.55 proves that both are nonidentity on
their domain by full R2 witnesses, and B.57 refutes equality
of their corrected common kernel with $\mathcal K$.
The assertion that every element of $\widehat{GT}\cap D$
is fixed by the two quintic maps remains open. The arithmetic
identity on $G_k$ is not promoted to that assertion on
all of $\widehat{GT}$.

The finite permutation and affine marking computations are also
checked exhaustively in \path{code/check_resolvent_character.py};
its output is retained. This finite certificate supplements the
restriction, quotient, and arithmetic arguments above.
